\section{Control de bajo nivel}
\label{sec:control}

\subsection{Arquitectura en cascada}

El control de actitud es una cascada de tres niveles, la misma de PX4, del
firmware ESP32 del proyecto y de la descripción que hace Larotta
\cite{larotta24} (p.~57) del Crazyflie (\cod{control/attitude.py}):
\begin{equation}
\underset{\scriptscriptstyle \text{actitud}}{(\phi,\theta,\psi)_{ref}}
\xrightarrow{\ \text{P}\ }
\underset{\scriptscriptstyle \text{tasas de Euler}}{(\dot\phi,\dot\theta,\dot\psi)_{ref}}
\xrightarrow{\ W\ }
\underset{\scriptscriptstyle \text{tasa}}{\vect\omega_{ref}}
\xrightarrow{\ \text{PID}\ }
\underset{\scriptscriptstyle \text{momento norm.}}{u\in[-1,1]^3}
\xrightarrow{\ \times\tau_{max}\ }
\vect\tau
\xrightarrow{\ \text{mezclador}\ }\sigma_{1..4}.
\label{eq:cascada}
\end{equation}
La cascada funciona porque las constantes de tiempo están separadas: el lazo
de tasa ve un integrador rápido (\cref{eq:Krate}) y el de ángulo ve el lazo de
tasa cerrado como un retardo corto. Powers, Mellinger y Kumar \cite{hb16}
(p.~317) lo justifican desde la linealización en hover: a ángulos pequeños
el control de actitud se desacopla del de posición y admite lazos anidados.

\paragraph{Lazo de ángulo.} Proporcional puro sobre el error de cada ángulo
de Euler (el de yaw envuelto a $(-\pi,\pi]$ para girar por el camino corto),
más una prealimentación de tasa de yaw opcional:
\begin{equation}
\dot\phi_{ref}=K_{\phi}(\phi_{ref}-\hat\phi),\quad
\dot\theta_{ref}=K_{\theta}(\theta_{ref}-\hat\theta),\quad
\dot\psi_{ref}=K_{\psi}\operatorname{wrap}(\psi_{ref}-\hat\psi)+\dot\psi_{ff},
\end{equation}
convertido a tasa del cuerpo con \eqref{eq:W} y acotado a $220^\circ$/s en
roll y pitch y $120^\circ$/s en yaw. No lleva integrador: el lazo de tasa
interno ya tiene uno, y un segundo integrador sobre una planta integradora
solo agrega sobreimpulso.

\paragraph{Lazo de tasa.} Un PID por eje (\cod{control/pid.py}) cuya salida
es el momento normalizado $u$; el momento físico es
$\tau=u\,\tau_{max}$ con $\tau_{max}$ de \eqref{eq:taumax}. Normalizar es la
convención de PX4 (\cod{MC\_ROLLRATE\_P} está en momento normalizado por
rad/s) y tiene una ventaja práctica: las ganancias de vehículos con
inercias que difieren en tres órdenes de magnitud quedan en el mismo rango
numérico.

\subsection{El PID discreto}

En cada periodo $\Delta t$ (\cod{PID.update}):
\begin{align}
e_k&=r_k-y_k,\nonumber\\
I_k&=\begin{cases}I_{k-1} & \text{si la salida está saturada y } e_k\,u_{k-1}>0\\
\sat_{\pm i_{lim}/k_i}\!\big(I_{k-1}+e_k\Delta t\big) & \text{en otro caso}\end{cases}\label{eq:pid}\\
d_k&=d_{k-1}+\alpha_d\Big(-\frac{y_k-y_{k-1}}{\Delta t}-d_{k-1}\Big),\qquad \alpha_d=\frac{\Delta t}{\tau_d+\Delta t},\ \tau_d=\frac{1}{2\pi f_d}\nonumber\\
u_k&=\sat_{\pm u_{lim}}\big(k_pe_k+k_iI_k+k_dd_k+k_{ff}r_k\big).\nonumber
\end{align}
Cada protección responde a una falla real. La derivada sobre la
\emph{medida}, y no sobre el error, evita el golpe derivativo cuando la
referencia salta. El filtro del derivativo ($f_d=40$~Hz nominales en el
lazo de tasa) evita que la vibración llegue amplificada a los motores. Se
discretiza como en Betaflight, $\alpha=\Delta t/(\tau+\Delta t)$, y esa
fórmula no da el corte nominal: a 500~Hz el corte efectivo es 32.4~Hz, con
unos 4.9~ms de retardo en baja frecuencia. Se conserva por compatibilidad con
las ganancias afinadas; la opción \cod{d\_filter\_exact} de la clase \cod{PID} usa
$\alpha=1-e^{-2\pi f_d\Delta t}$, cuyo corte coincide con el nominal. Cuál de
las dos usa el firmware del dron está por verificar. La integración
condicional congela el integrador cuando la salida está saturada y el error
empuja hacia el mismo lado (\emph{anti-windup}), y el límite absoluto
$|k_iI|\le i_{lim}=0.3$ acota cuánto momento puede acumular. En el lazo de
tasa $u_{lim}=1$.

\subsection{Sintonía SIMC a partir del modelo}

Las ganancias por defecto no están escritas a mano: \cod{control/tuning.py}
las calcula con las reglas SIMC de Skogestad \cite{skogestad03} a partir de
los parámetros físicos, de modo que cambiar de vehículo cambia las
ganancias. SIMC tiene un solo parámetro libre, la constante de tiempo de
lazo cerrado deseada $\tau_c$, con interpretación física directa; es más
fácil de defender ante evaluadores que una sintonía por optimización de caja
negra, y produce márgenes de robustez razonables por construcción
(Ziegler--Nichols, en comparación, da sobreimpulsos del orden del 25~\%
\cite{astrom04}).

\paragraph{Modelo de diseño del lazo de tasa.} Linealizando en hover
(\cref{sec:linealizacion}) cada eje es un integrador, y los retardos de la
cadena se agrupan con la ``regla de la mitad'' de SIMC (la mitad de la
constante de tiempo despreciada se suma al retardo):
\begin{equation}
\frac{\omega(s)}{u(s)}=\frac{K\,e^{-\theta s}}{s},\qquad
K=\frac{\tau_{max}}{J},\qquad
\theta=\underset{\scriptscriptstyle 2\ \text{ms}}{\tau_{IMU}}+\underset{\scriptscriptstyle 1\ \text{ms}}{\Delta t/2}+\underset{\scriptscriptstyle 3\ \text{ms}}{\tau_{filtro}}+\underset{\scriptscriptstyle 10\ \text{ms}}{T_m/2}=\val{theta}~\text{ms}.
\label{eq:simc_planta}
\end{equation}
Esa es la cadena \emph{nominal} con la que se afinaron y documentaron las
ganancias por defecto: 2~ms de IMU y 3~ms fijos de filtro, el retardo del
Butterworth de 80~Hz de las primeras versiones (\val{gdBWo}~ms de retardo de
grupo a 10~Hz, \cref{tab:filtros}). Con el filtro actual, \val{filtroDef}, el
retardo de grupo a 10~Hz es \val{defGd10}~ms. \cod{chain\_delay} también sabe
leer la cadena \emph{configurada} (\cod{Simulation.attitude\_chain}): retardo
de la IMU de \cod{SensorConfig}, retardo de grupo de las etapas de filtro de
cada eje evaluado en el cruce estimado (que depende de $\theta$ y se
resuelve por punto fijo), retardo de cómputo y retardo del ESC. Con la
configuración por defecto da $\theta=16.6$~ms en lugar de 16.0, y un $K_p$ de
tasa 3.6~\% menor. La simulación solo usa la cadena configurada si se pide
(\cod{SimConfig.gains\_from\_chain}), para que las ganancias no cambien en
silencio durante un barrido de filtros, donde el filtro varía con las
ganancias fijas. Dos cosas quedan fuera de $\theta$ por diseño: el filtro del
derivativo, que solo afecta a esa rama, y la forma exacta del rezago del
motor. La regla de la mitad reparte $T_m/2=10$~ms, cuando el motor de primer
orden aporta en la frecuencia de cruce más fase que esa. Las ganancias SIMC
son, por eso, un punto de partida; los márgenes reales del lazo se discuten
en la \cref{sec:integracion}. Para el PAE-X450,
$K_x=\val{KR}$, $K_y=\val{KP}$ y $K_z=\val{KY}$~rad/s$^2$.

\paragraph{Reglas.} Para un integrador con retardo, SIMC da un PI
\begin{equation}
K_c=\frac{1}{K(\tau_c+\theta)},\qquad T_i=4(\tau_c+\theta),\qquad k_i=\frac{K_c}{T_i},
\qquad \tau_c=\theta .
\label{eq:simc}
\end{equation}
La elección $\tau_c=\theta$ es la recomendación ``ajustada'' de Skogestad
(margen de ganancia $\approx3$, margen de fase $\approx47^\circ$). El código
agrega un derivativo en roll y pitch, $k_d=K_c\,\theta$ (es decir,
$T_d=\theta$), porque la planta real tiene el rezago del motor que la regla
de la mitad solo reparte y que el derivativo compensa con adelanto de fase;
en yaw $k_d=0$. El factor no viene de SIMC sino de los barridos automáticos
(\cref{sec:barridos}): la primera versión usaba $k_d=K_c\theta/3$, y frente a
un doblete de \emph{tasa} ese derivativo dejaba un 20~\% de sobreimpulso;
con $K_c\theta$ y el filtro de 60~Hz baja a cerca del 6~\%, a costa de un 7~\%
más de actividad de alta frecuencia en los motores. El lazo de
ángulo es proporcional, diseñado sobre $1/s$ con el lazo de tasa cerrado
visto como un retardo $\tau_{cl}=\tau_c+\theta$ y una constante de lazo
cerrado $3\tau_{cl}$:
\begin{equation}
K_{ang}=\frac{1}{3\tau_{cl}+\tau_{cl}}=\frac{1}{4\tau_{cl}}.
\label{eq:kang}
\end{equation}

\paragraph{Con los números del PAE-X450.} $\theta=16$~ms y
$\tau_c+\theta=32$~ms, de modo que
\begin{align*}
\text{roll:}&\quad K_c=\frac{1}{\val{KR}\cdot0.032}=\val{kpRateRoll},\quad T_i=0.128~\text{s},\quad k_i=\val{kiRateRoll},\quad k_d=\val{kdRoll};\\
\text{pitch:}&\quad K_c=\val{kpRatePitch},\quad k_i=\val{kiRatePitch},\quad k_d=\val{kdPitch};\\
\text{yaw:}&\quad K_c=\val{kpRateYaw},\quad k_i=\val{kiRateYaw},\quad k_d=0;\\
\text{ángulo:}&\quad K_{ang}=\frac{1}{4\cdot0.032}=\val{kpAng}~\text{s}^{-1}\ \text{(los tres ejes)}.
\end{align*}
Como comprobación, la frecuencia de cruce del lazo de tasa de roll es
$K_cK=\val{wc}$~rad/s $=1/(\tau_c+\theta)$, como prescribe SIMC. En unidades
físicas, $K_c\tau_{max}=\val{KpNm}$~N·m por rad/s. El $K_c$ de yaw es 16 veces
mayor que el de roll porque su planta es 16 veces más lenta
($K_z\ll K_x$): las ganancias normalizadas absorben la diferencia de
autoridad.

\paragraph{Respuesta medida.} En el simulador completo (estimador ideal,
sensores sin ruido), un escalón de roll de $10^\circ$ tiene un tiempo de
subida del 10 al 90~\% de \val{stepRise}~ms, sobreimpulso
\val{stepOS}~\% y error final \val{stepErr}$^\circ$. La subida es más rápida
que la que sugiere el objetivo de diseño (una constante de lazo cerrado de
$3\tau_{cl}=96$~ms daría $2.2\times96\approx210$~ms de subida), señal de que la regla de la
mitad sobrestima el retardo efectivo: el diseño queda del lado conservador.
Con un estimador real en el lazo, la identificación de la
\cref{sec:posicion} da un primer orden con retardo de
$\val{tauAttF}~\text{ms}+\val{thAttF}~\text{ms}$ para el lazo de actitud
cerrado.

\subsection{Lazo vertical}

En modo actitud, la altura se controla con una cascada posición--velocidad
que corre en el lazo externo (100~Hz) y entrega una aceleración vertical
deseada hacia arriba (\cod{VerticalController}):
\begin{equation}
\begin{gathered}
\dot h_{ref}=\sat_{\pm2.5}\big(K_h(h_{ref}-\hat h)\big),\qquad
a_{up}=\text{PI}\big(\dot h_{ref}-\hat{\dot h}\big)\in[-6,6]~\text{m/s}^2,\\
T=\sat_{[0,\,0.95T_{max}]}\frac{m\,(g+a_{up})}{\max(\cos\hat\phi\cos\hat\theta,\ 0.5)} .
\end{gathered}
\label{eq:vertical}
\end{equation}
La división por $\cos\phi\cos\theta$ compensa la pérdida de empuje vertical
al inclinarse (con un tope en $60^\circ$ de inclinación equivalente). Con la
compensación de masa, la planta entre $a_{up}$ y $\dot h$ es un integrador
con ganancia 1; las ganancias salen de las mismas reglas con un retardo
supuesto de 50~ms y $\tau_c=2.5\theta$: $K_h=\val{vKpPos}$~s$^{-1}$,
$k_p=\val{vKpVel}$~s$^{-1}$ y $k_i=\val{vKiVel}$~s$^{-2}$.

\subsection{De la aceleración deseada a empuje y actitud: planitud diferencial}
\label{sec:planitud}

Los cinco controladores de posición (\cref{sec:posicion}) entregan una
aceleración deseada $\vect a_{des}$ en NED, y una sola función
(\cod{acceleration\_to\_attitude}) la convierte en empuje y actitud de
referencia. Mellinger y Kumar \cite{mellinger11} mostraron que el
cuadricóptero es diferencialmente plano con salidas planas $(\vect p,\psi)$:
la actitud y el empuje quedan determinados por la aceleración y el rumbo.
De \eqref{eq:newton} sin arrastre, la fuerza que debe producir el empuje es
\begin{equation}
\vect F=m\,(\vect a_{des}-g\eu),\qquad F_z\le-0.1\,mg,
\qquad \|\vect F_{xy}\|\le|F_z|\tan\theta_{max},
\label{eq:flatF}
\end{equation}
donde la primera cota impide pedir empuje hacia abajo y la segunda limita
la inclinación a $\theta_{max}=\val{tiltMax}^\circ$
(\cod{SimConfig.tilt\_max\_deg}). El simulador pasa ese mismo límite a
\cod{acceleration\_to\_attitude} para todos los controladores, de modo que
ninguno puede inclinar más que otro. Como el empuje actúa a lo largo de $-z_B$:
\begin{equation}
\vect z_B=-\frac{\vect F}{\|\vect F\|},\quad
\vect x_c=(\cos\psi_{ref},\sin\psi_{ref},0),\quad
\vect y_B=\frac{\vect z_B\times\vect x_c}{\|\vect z_B\times\vect x_c\|},\quad
\vect x_B=\vect y_B\times\vect z_B,
\label{eq:flatR}
\end{equation}
y de $\R=[\vect x_B\ \vect y_B\ \vect z_B]$ se leen
$\theta=-\arcsin x_{B,z}$, $\phi=\operatorname{atan2}(y_{B,z},z_{B,z})$,
$\psi=\operatorname{atan2}(x_{B,y},x_{B,x})$ y $T=\|\vect F\|$. Con el
PAE-X450, pedir 1~m/s$^2$ hacia el norte da $\theta=\val{flatPitch}^\circ$
(morro abajo) y $T=\val{flatT}$~N; 2~m/s$^2$ hacia el este da
$\phi=\val{flatRoll}^\circ$. Ambos coinciden con
$\arctan(a/g)$ y $m\sqrt{g^2+a^2}$.

Usar una conversión común es lo que hace justa la comparación: los cinco
controladores ven la misma planta, los mismos sensores, el mismo estimador y
el mismo bajo nivel, y solo difieren en cómo calculan $\vect a_{des}$.
